% Einheit 5 von 5 – Das Integral als Flächenbilanz (bank/flaecheninhalt-durch-integration/e5.jsonl, zusammenbau v0.1)
%% TODO Zeitmarke Einheit 5: Marken-Zeile nicht lesbar
%% TODO Zweigzeile Teil 1: was hier gelernt wird (Satz in Schülersprache aus dem Einheitstitel) – Einheit 5
\einheitenkopf[e5]{Einheit 5 von 5 · Das Integral als Flächenbilanz}
\zweigzeile{Abitur GK}
\setcounter{aufgabe}{16}

%% TODO Titel: Ich-kann-Satz fehlt in der Bank (Platzhalter: Kettenname) – Nr. 17
%% TODO Anweisung: ein Satz über der ersten Teilaufgabe fehlt in der Bank – Nr. 17
\begin{aufgabe}{Flächenbilanz}
%% TODO \rechenplatz{n}: Schrittzahl des Grundfalls fehlt in der Bank (nur bei fester Schreibform, 2.3 b)
\begin{teile}
\teil Integral von $0$ bis $2$ über $f(x)$, Graph im Bild – Vorzeichen des Werts? Ist der Wert der Flächeninhalt? \kreuz{Wert positiv}\\ \kreuz{Wert negativ}\\ \kreuz{Wert null}

\begin{ksys}[xmin=-3,xmax=5,ymin=-4,ymax=4,xstep=1,ystep=1,ablesen]\funktion{\x^2-4}{f}\end{ksys}
\teil Integral von $0$ bis $4$ über $f(x)$ – bestimme den Wert durch Kästchenzählen am Bild.

\begin{ksys}[xmin=-1,xmax=5,ymin=-4,ymax=4,xstep=1,ystep=1,ablesen]\funktion{0.5*\x+1}{f}\end{ksys}
\teil Begründe am Bild, dass das Integral von $-2$ bis $2$ über $f(x)$ null ist.

\begin{ksys}[xmin=-3,xmax=5,ymin=-4,ymax=4,xstep=1,ystep=1,ablesen]\funktion{0.5*\x^3-2*\x}{f}\end{ksys}
\teil $f(x) = 2x^3 - 5x + 3$ – Integral von $-2$ bis $2$ ohne Stammfunktion?
\teil Der Graph von $f$ liegt zwischen $0$ und $3$ über, zwischen $3$ und $5$ unter der x-Achse. Ist das Integral von $0$ bis $5$ kleiner, gleich oder größer als das von $0$ bis $3$? \kreuz{kleiner}\\ \kreuz{gleich}\\ \kreuz{größer}
\teil $f(x) = x^2 + 1$ – nähere das Integral von $0$ bis $2$ durch das Trapez unter der Sehne an. Über- oder Unterschätzung?
\teil $f(x) = 3x^2$ – Mittelwert der Funktionswerte auf dem Intervall von $0$ bis $2$?
\teil Für $k \in \mathbb{R}$ ist $f_k(x) = (x - k)^2 \cdot e^{x}$. Beurteile ohne Berechnung eines Integrals die Aussage: Für jeden Wert von $k$ gibt das Integral von $-2$ bis $2$ über $f_k(x)$ den Inhalt der Fläche zwischen Graph und x-Achse auf diesem Intervall an. \hfill (Abitur 2025 LK)
\end{teile}
\end{aufgabe}

%% TODO Titel: Ich-kann-Satz fehlt in der Bank (Platzhalter: Kettenname) – Nr. 18
%% TODO Anweisung: ein Satz über der ersten Teilaufgabe fehlt in der Bank – Nr. 18
\begin{aufgabe}{Flächenbilanz – Fehler finden, Begründen, Anwendung, Darstellungswechsel}
\begin{teile}
\teil Fläche zwischen dem Graphen von $f(x) = x^3$ und der x-Achse von $x = -2$ bis $x = 1$ – wo ist der Fehler? \rechnung{F(1) - F(-2) &= \tfrac{1}{4} - 4 = -\tfrac{15}{4} \\ A &= \tfrac{15}{4}}
\teil Warum ist das Integral einer ungeraden Funktion über einem zu null symmetrischen Intervall null?
\teil Eine Solaranlage liefert die Leistung $p(t) = -0{,}25t^2 + 3t$ (in kW, $t$ in Stunden ab 6 Uhr, $0 \le t \le 12$). Mittlere Leistung – und was bedeutet sie?
\teil Zeichne im Bild ein Flächenstück zwischen dem Graphen von $f(x) = 0{,}5x + 1$, der x-Achse, der y-Achse und einer Parallelen zur y-Achse mit dem Inhalt $3$ ein und gib einen Term für seinen Inhalt an.

\begin{ksys}[xmin=-1,xmax=5,ymin=-1,ymax=5,ablesen]\funktion{0.5*\x+1}{f}\end{ksys}
\end{teile}
\end{aufgabe}
